% ====== 高中物理笔记 · 公共导言区（所有章节共用，章节文件里只能使用这里定义的宏/环境）======
\usepackage{amsmath,amssymb,mathtools,bm}
\usepackage{graphicx,xcolor,enumitem,booktabs,multirow,makecell,array,colortbl,diagbox,extarrows}
\usepackage[normalem]{ulem}
\usepackage{tikz}
\usetikzlibrary{arrows.meta,calc,positioning,shapes.geometric,decorations.pathmorphing}
\usepackage{tcolorbox}
\tcbuselibrary{skins,breakable}
% XeLaTeX：可跨页框被拆开时，盒子里残留的颜色 push 会打乱 pgf 的颜色栈技巧，在栈里留下白色，导致后文正文变成白字。
% 去掉盒子内部的颜色设置即可（正文、标题颜色在绘制时另行设置，不受影响）。pdfTeX/LuaTeX 有独立颜色栈，不需要。
\ifdefined\XeTeXrevision
\makeatletter
\patchcmd{\tcb@vbox}{\color{#3}}{}{}{\PackageWarning{preamble}{could not patch tcb@vbox}}
\makeatother
\fi
\usepackage{geometry}
\geometry{a4paper,left=2.2cm,right=2.2cm,top=2.4cm,bottom=2.4cm}
\usepackage{fancyhdr}
\usepackage[colorlinks=true,linkcolor=blue!50!black,urlcolor=blue!50!black]{hyperref}
\setlist{nosep,leftmargin=*,itemsep=2pt}
\setlist[itemize,1]{label=\textbullet}
\setlist[itemize,2]{label=\textendash}
\setlist[itemize,3]{label=\ensuremath{\ast}}

% ---- 红笔批注 / 辨认不清 / 原稿划去 ----
\newcommand{\red}[1]{\textcolor{red!75!black}{#1}}
\newcommand{\unclear}[1]{\textcolor{orange!80!black}{[#1\,\textbf{?}]}}
\newcommand{\struck}[1]{\sout{#1}}

% ---- 常用数学/物理记号 ----
\newcommand{\dd}{\mathrm{d}}
\newcommand{\nuc}[3]{{}^{#1}_{#2}\mathrm{#3}}          % 核素：\nuc{238}{92}{U}
\newcommand{\unit}[1]{\,\mathrm{#1}}                    % 单位：5\unit{m/s}
\newcommand{\ee}{\mathrm{e}}
% fontspec 下大写希腊字母放进 \mathrm（如 \unit{\Omega/m}）会落到文本字体的空位而整个消失；固定到 OT1 字体族即可
% （fontspec 在 \begin{document} 时才声明希腊字母，所以这里也放在 \AtBeginDocument 里、在它之后执行）
\DeclareSymbolFont{pnupgreek}{OT1}{cmr}{m}{n}
\AtBeginDocument{%
  \DeclareMathSymbol{\Gamma}{\mathord}{pnupgreek}{"00}%
  \DeclareMathSymbol{\Delta}{\mathord}{pnupgreek}{"01}%
  \DeclareMathSymbol{\Theta}{\mathord}{pnupgreek}{"02}%
  \DeclareMathSymbol{\Lambda}{\mathord}{pnupgreek}{"03}%
  \DeclareMathSymbol{\Xi}{\mathord}{pnupgreek}{"04}%
  \DeclareMathSymbol{\Pi}{\mathord}{pnupgreek}{"05}%
  \DeclareMathSymbol{\Sigma}{\mathord}{pnupgreek}{"06}%
  \DeclareMathSymbol{\Upsilon}{\mathord}{pnupgreek}{"07}%
  \DeclareMathSymbol{\Phi}{\mathord}{pnupgreek}{"08}%
  \DeclareMathSymbol{\Psi}{\mathord}{pnupgreek}{"09}%
  \DeclareMathSymbol{\Omega}{\mathord}{pnupgreek}{"0A}%
}

% ---- 盒子环境 ----
\newtcolorbox{keybox}[1][]{breakable,enhanced,colback=yellow!6,colframe=orange!70!black,
  boxrule=0.6pt,arc=2pt,left=6pt,right=6pt,top=4pt,bottom=4pt,title={#1},fonttitle=\bfseries,
  attach boxed title to top left={xshift=6pt,yshift=-3mm},boxed title style={colback=orange!70!black,arc=2pt}}
\newtcolorbox{notebox}[1][注意]{breakable,enhanced,colback=red!4,colframe=red!60!black,
  boxrule=0.5pt,arc=2pt,left=6pt,right=6pt,top=3pt,bottom=3pt,title={#1},fonttitle=\bfseries\small}
\newtcolorbox{problem}[1][]{breakable,enhanced,colback=blue!3,colframe=blue!55!black,
  boxrule=0.6pt,arc=2pt,left=6pt,right=6pt,top=4pt,bottom=4pt,title={例题\ #1},fonttitle=\bfseries}
\newenvironment{solution}{\par\smallskip\noindent\textbf{解：}\ }{\par\smallskip}

% ---- 手写图片（从原稿裁剪）：\notefig[宽度占行宽比例]{figures/文件名.png} ----
\newcommand{\notefig}[2][0.6]{\begin{center}\includegraphics[width=#1\linewidth,height=0.45\textheight,keepaspectratio]{#2}\end{center}}

% ---- 文本模式下的特殊符号交给中文字体（①②… ★☆○●■□△）----
\xeCJKsetcharclass{"2460}{"2473}{1}
\xeCJKsetcharclass{"25A0}{"25EF}{1}
\xeCJKsetcharclass{"2605}{"2606}{1}
